\documentclass[aps, reprint, twocolumn, superscriptaddress, floatfix]{revtex4-2}

% --- Writing Packages ---
\usepackage{amsmath, amssymb}
\let\Bbbk\relax % Avoid conflict with newtxmath
\usepackage{newtxmath, newtxtext} % Load after amsmath
\usepackage{graphicx}
\usepackage{tikz}
\usepackage{tikz-3dplot}
\usetikzlibrary{arrows.meta, angles, quotes, topaths, shadings, decorations.markings}
\usepackage{quantikz}      % Quantum circuit diagrams
\usepackage{booktabs}
\usepackage{tikzorbital}
\usepackage{float}
\usepackage{hyperref}
\usepackage{braket}      % Dirac notation: \ket{}, \bra{}, \braket{}
\usepackage{physics}     % \tr, \norm, \abs, etc.
\usepackage{mathtools}   % \coloneqq, dcases, etc.
\hypersetup{
    colorlinks=true,
    linkcolor=blue,
    citecolor=blue,
    urlcolor=blue
}

% --- Custom Environments ---
% Postulate / Theorem statement block (replaces \begin{quote})
\newenvironment{statement}{%
  \par\addvspace{5pt}%
  \noindent\rule{\linewidth}{0.4pt}\nopagebreak\par\vspace{2pt}%
  \noindent\begingroup\small\itshape}{%
  \endgroup\par\vspace{2pt}\nopagebreak%
  \noindent\rule{\linewidth}{0.4pt}\par\addvspace{5pt}}

% Proof-sketch environment (avoids amsthm/RevTeX conflict)
\newenvironment{proofsketch}[1][Sketch of proof]{%
  \par\smallskip\noindent\textit{#1.}\enspace}{%
  \hspace*{\fill}$\lozenge$\par\smallskip}

\usepackage{style} % Style definitions only

\begin{document}

\preprint{STA4121. QUANTUM MACHINE LEARNING}

\title{An Introduction to Quantum Machine Learning}

\author{Jeongbin Jo}
\email{jeongbin033@yonsei.ac.kr}
\affiliation{Department of Physics, Yonsei University, Seoul, 03722, Republic of Korea}

\date{\today}

\begin{abstract}
Quantum machine learning (QML) sits at the intersection of quantum information
science and modern artificial intelligence, offering a promising route toward
computational advantages that are inaccessible by classical means alone.
This report provides a self-contained introduction to the theoretical foundations
necessary to understand and critically evaluate QML proposals.  Starting from the
postulates of quantum mechanics---state, dynamics, and measurement---we develop
the quantum circuit model with single- and multi-qubit gates, establish the
universality of quantum gate sets, and examine the no-cloning theorem and
elementary quantum communication protocols.  We then discuss reversible
computation, Landauer's principle, the Hadamard and swap tests with complete
circuit derivations, and the basics of quantum computational complexity,
culminating in an overview of landmark QML algorithms (HHL, QPCA, QSVM) and the
central open questions surrounding quantum advantage in machine learning tasks.
Throughout, formal derivations are presented at a level of detail beyond standard
textbook treatments, with the aim of equipping the reader with both intuition and
rigorous mathematical foundations.
\end{abstract}

\maketitle
\thispagestyle{fancytitle}
% --- Watermark (Seal) ---
\addwatermark

\onecolumngrid  % 여기서부터 1단 시작
\vspace{-1cm}
\tableofcontents
\vspace{1cm}    % 목차와 본문 사이 간격 조정
\twocolumngrid  % 여기서부터 다시 2단 시작

% ============================================================
\section{Introduction}
\label{sec:introduction}
% ============================================================

\subsection{Motivation: The Age of Quantum Computing}
\label{ssec:motivation}

For over half a century, classical computing has followed Moore's Law: the
number of transistors on a microprocessor doubles approximately every two years,
delivering exponential growth in computational power at constant cost
\cite{preskill2018quantum}.  However, as transistor feature sizes approach
atomic dimensions, quantum mechanical effects such as tunneling and thermal
fluctuations undermine deterministic switching.  This physical barrier signals
the end of the classical scaling era and motivates the search for radically
different computational paradigms.

Quantum computing exploits the principles of quantum mechanics---superposition,
entanglement, and interference---to represent and process information in ways
that classical computers cannot efficiently simulate.  A classical register of
$n$ bits can store exactly one of $2^n$ binary strings at any given time.  By
contrast, an $n$-qubit quantum register can exist in a superposition of all
$2^n$ computational basis states simultaneously (see Postulate~I in 
Section~\ref{ssec:state} for a formal definition).
The amplitudes $\{\alpha_x\}$ encode correlations across an exponentially large
state space, and unitary operations manipulate all $2^n$ amplitudes in
parallel---a wave-like processing step.  Measurement then collapses this
superposition, yielding a bit-string $x$ with probability $|\alpha_x|^2$---a
particle-like readout step.  This \emph{analog-digital duality} underlies the
power of quantum computing \cite{nielsen2010quantum}.

\subsection{Classical vs.\ Quantum Information}
\label{ssec:classical-vs-quantum}

Classical information theory studies the representation, compression, and
transmission of information encoded in discrete, deterministic symbols
\cite{wilde2017quantum}.  A probabilistic extension replaces each bit state by a
probability distribution $\mathbf{p} = (p_0, p_1)^\top$, $p_0 + p_1 = 1$,
$p_i \ge 0$.  Operations on probabilistic bits are described by stochastic
matrices---matrices whose columns are probability vectors.

Quantum information theory is a mathematical generalization of classical
probabilistic information.  Specifically, whereas classical transitions are
modeled by a \emph{Markov chain} where a stochastic matrix $M$ acts on a
probability vector $\mathbf{x}$, quantum transitions are governed by
\emph{unitary evolution} where a unitary matrix $U$ acts on a ket state 
$\ket{\psi}$. The key correspondences are summarized below:
\begin{center}
\resizebox{\columnwidth}{!}{%
\begin{tabular}{lll}
\toprule
 & Classical (Markov) & Quantum (Unitary) \\
\midrule
State & Prob.\ vector $\mathbf{x} \in \mathbb{R}_{+}^{d}$ & Ket $\ket{\psi} \in \mathbb{C}^{d}$ \\
Normalization & $L_1$ norm ($\sum x_i = 1$) & $L_2$ norm ($\sum |\alpha_i|^2 = 1$) \\
Evolution & Stochastic $M$ ($\sum_i M_{ij} = 1$) & Unitary $U$ ($U^\dagger U = \mathbf{I}$) \\
Observation & Sampling from $\mathbf{x}$ & Born-rule measurement \\
\bottomrule
\end{tabular}}
\end{center}
This structural analogy reveals that quantum mechanics is, in a precise sense, a
\emph{complex-amplitude} generalization of probability theory \cite{watrous2018theory}.
The crucial difference is that complex amplitudes can interfere constructively and
destructively---classical probabilities cannot.  Interference is the key
resource exploited by quantum algorithms to suppress wrong answers and amplify
correct ones \cite{nielsen2010quantum}.

\subsection{Machine Learning as an Optimization Problem}
\label{ssec:ml-opt}

Machine Learning (ML) aims to enable computers to learn underlying properties of
data to perform tasks without explicit programming. Mathematically, ML represents
an optimization problem where we search for a model $f(z,x)$ parameterized by $z$
that minimizes an expected loss $L$ over a data distribution $\mathcal{D}$
\cite{schuld2021machine}:
\begin{equation}
  \min_{z \in \Omega} \mathbb{E}_{x \sim \mathcal{S}} \bigl\{ L[f(z, x)] \bigr\},
\end{equation}
where $\mathcal{S} = \{x_i\}_{i=1}^m \sim \mathcal{D}$ is the available sample.
Typical objective functions can be derived from the principle of 
\emph{Maximum Likelihood Estimation} (MLE). For a dataset $Y = \{y_1, \dots, y_m\}$,
the log-likelihood $L(\theta)$ is maximized:
\begin{equation}
  \max_{\theta} L(\theta) = \sum_{i=1}^m \log(p(y_i | x_i, \theta)).
\end{equation}
Under the assumption of i.i.d.\ Gaussian noise $\epsilon \sim \mathcal{N}(0, \sigma^2)$
in observations $y_i = f(x_i, \theta) + \epsilon$, MLE is equivalent to minimizing
the sum of squared errors (SSE/MSE):
\begin{equation}
  L(\theta) \propto -\sum_{i=1}^m (y_i - f(x_i, \theta))^2 + \text{const}.
\end{equation}
Similarly, for classification tasks with Bernoulli labels $y_i \in \{0, 1\}$, 
MLE yields the cross-entropy loss function:
\begin{equation}
  L(\theta) = \sum_{i=1}^m \bigl[ y_i \log \theta_i + (1-y_i) \log (1-\theta_i) \bigr].
\end{equation}

\subsection{Scope and Organization of This Report}
\label{ssec:scope}

This report covers the foundational material of STA4121 (Lectures 1--6) at a
level of mathematical detail intended to complement and, in key derivations,
surpass standard textbook treatments.  The organization is as follows.
Section~\ref{sec:introduction} introduces the motivation and the machine 
learning optimization framework. Section~\ref{sec:math} reviews prerequisites 
including Dirac notation, computational complexity, and classical information theory.
Section~\ref{sec:postulates} develops the three postulates of quantum mechanics.
Section~\ref{sec:multiqubit} treats entanglement and density matrices.
Section~\ref{sec:circuits} introduces the quantum circuit model and 
elementary protocols including Superdense Coding and Teleportation.
Section~\ref{sec:reversibility} discusses reversible computation, 
Landauer's principle, and the Hadamard/swap tests.
Section~\ref{sec:qml} gives an overview of key QML algorithms.
Section~\ref{sec:conclusion} concludes.

% ============================================================
\section{Mathematical Preliminaries}
\label{sec:math}
% ============================================================

\subsection{Dirac Notation and Hilbert Space}
\label{ssec:dirac}

All quantum states live in a complex Hilbert space $\mathcal{H}$, a complete
inner-product space over $\mathbb{C}$.  Following Dirac's notation:
\begin{itemize}
  \item A \emph{ket} $\ket{\psi} \in \mathcal{H}$ is a column vector representing
        a quantum state.
  \item A \emph{bra} $\bra{\psi} = \ket{\psi}^\dagger$ is the conjugate-transpose
        row vector.
  \item The \emph{inner product} $\braket{\phi|\psi} \in \mathbb{C}$ satisfies
        linearity in the second argument and $\braket{\psi|\psi} \ge 0$, with
        equality iff $\ket{\psi} = 0$.
  \item The \emph{outer product} $\ket{\psi}\bra{\phi}$ is a linear operator on
        $\mathcal{H}$.
\end{itemize}
For an $n$-qubit system, $\mathcal{H} = \mathbb{C}^{2^n}$ with the standard
inner product $\braket{\phi|\psi} = \sum_x \phi_x^* \psi_x$.  The
\emph{computational basis} $\{\ket{x}\}_{x \in \{0,1\}^n}$ is the set of
standard basis vectors, which forms a complete orthonormal set satisfying
$\braket{x|y} = \delta_{xy}$ and $\sum_x \ket{x}\bra{x} = \mathbf{I}$.

\subsection{Asymptotic Complexity Notation}
\label{ssec:complexity-notation}

In the analysis of both classical and quantum algorithms, it is essential to characterize how the computational cost (time or space) scales with the input size $N$ as $N \to \infty$ \cite{bernstein1997quantum}.
\begin{enumerate}
    \item \textbf{Big-O ($O$)}: $f(N) = O(g(N))$ if there exist constants $c, N_0 > 0$ such that $f(N) \le c \cdot g(N)$ for all $N \ge N_0$. This provides an \emph{upper bound} on growth.
    \item \textbf{Big-Omega ($\Omega$)}: $f(N) = \Omega(g(N))$ if $f(N) \ge c \cdot g(N)$ for all $N \ge N_0$. This provides a \emph{lower bound}.
    \item \textbf{Big-Theta ($\Theta$)}: $f(N) = \Theta(g(N))$ if $f(N) = O(g(N))$ and $f(N) = \Omega(g(N))$. This represents a \emph{tight bound}.
    \item \textbf{Little-o ($o$)} and \textbf{Little-omega ($\omega$)}: Used for strict upper and lower bounds, respectively.
\end{enumerate}
An algorithm is generally considered \emph{efficient} if its complexity is polynomial in $N$, i.e., $O(N^k)$ for some constant $k$. Key examples include:
\begin{itemize}
    \item Inner product $\mathbf{x}^\top\mathbf{y}$ costs $O(N)$.
    \item Matrix-vector product $A\mathbf{x}$ costs $O(MN)$ for $A\in\mathbb{R}^{M\times N}$.
    \item Classical matrix inversion costs $O(N^3)$, while the HHL algorithm offers an exponential speedup to $O(\text{poly}\log N)$ for sparse, well-conditioned systems \cite{harrow2009quantum}.
\end{itemize}

\subsection{Classical Information: Probabilistic States and Stochastic Matrices}
\label{ssec:classical-info}

Let $\mathcal{S} = \{0,1\}$ denote the classical state set of a bit $X$.
A \emph{probabilistic state} of $X$ is a column vector
$\mathbf{p} = (p_0, p_1)^\top \in \mathbb{R}^2$ where $p_a = \Pr(X = a)$, so
$p_0, p_1 \ge 0$ and $p_0 + p_1 = 1$.

A \emph{deterministic operation} mapping $a \mapsto f(a)$ is represented by the
matrix $M$ satisfying $M\ket{a} = \ket{f(a)}$ for all $a \in \mathcal{S}$.
A \emph{probabilistic operation} introducing randomness is represented by a
\emph{stochastic matrix}: a matrix $M$ satisfying $M_{ij} \ge 0$ and $\sum_i
M_{ij} = 1$ for all $j$.  Stochastic matrices are exactly those that map
probability vectors to probability vectors.  Composition of operations:
\begin{equation}
  \mathbf{p} \xrightarrow{M_1} M_1 \mathbf{p} \xrightarrow{M_2}
  M_2 M_1 \mathbf{p},
\end{equation}
and the product $M_2 M_1$ is again stochastic.


% ============================================================
\section{Postulates of Quantum Mechanics}
\label{sec:postulates}
% ============================================================

\subsection{Postulate I: State}
\label{ssec:state}

\begin{statement}
\textbf{Postulate I (State).}
Any isolated physical system is completely described by a unit vector
$\ket{\psi}$ in a complex Hilbert space $\mathcal{H}$, called the state space.
For an $n$-qubit system $\mathcal{H}=\mathbb{C}^{2^n}$ and $\braket{\psi|\psi}=1$.
\end{statement}

For a single qubit, $\mathcal{H} = \mathbb{C}^2$ and the most general state is
\begin{equation}
  \ket{\psi} = \alpha\ket{0} + \beta\ket{1}, \quad
  \alpha,\beta \in \mathbb{C}, \quad
  |\alpha|^2 + |\beta|^2 = 1.
  \label{eq:qubit}
\end{equation}

\paragraph{Bloch sphere parameterization.}
Writing $\alpha = e^{i\gamma}\cos(\theta/2)$ and
$\beta = e^{i(\gamma+\phi)}\sin(\theta/2)$, the global phase $e^{i\gamma}$ is
physically unobservable (it cancels in every Born-rule probability
$p(m)=|\bra{m}\ket{\psi}|^2$).  Absorbing it, the canonical form is
\begin{equation}
  \ket{\psi} = \cos\!\tfrac{\theta}{2}\ket{0} + e^{i\phi}\sin\!\tfrac{\theta}{2}\ket{1},
  \quad \theta \in [0,\pi],\;\phi \in [0,2\pi).
  \label{eq:bloch}
\end{equation}
The pair $(\theta,\phi)$ defines a point on the unit sphere---the
\emph{Bloch sphere}---giving a bijective correspondence between physically
distinct pure single-qubit states and $S^2$.  Antipodal points correspond to
orthogonal states: the poles $\ket{0}$ and $\ket{1}$ are mutually orthogonal.

\begin{proofsketch}
We verify Eq.~\eqref{eq:bloch} represents all and only unit vectors (up to global phase).
Any $(\alpha,\beta) \in \mathbb{C}^2$ with $|\alpha|^2+|\beta|^2=1$ can be written
as $\alpha = r_0 e^{i\phi_0}$, $\beta = r_1 e^{i\phi_1}$ with $r_0^2+r_1^2=1$, $r_i\ge 0$.
Setting $r_0=\cos(\theta/2)$, $r_1=\sin(\theta/2)$ (unique for $\theta\in[0,\pi]$),
and factoring out $e^{i\phi_0}$ as the global phase gives $\phi = \phi_1-\phi_0$.
The parameterization is surjective onto the sphere and two-to-one only at the poles
($\theta=0,\pi$), where $\phi$ is irrelevant---consistent with $\ket{0},\ket{1}$
being pole states. $\lozenge$
\end{proofsketch}

\subsection{Postulate II: Dynamics}
\label{ssec:dynamics}

\begin{statement}
\textbf{Postulate II (Dynamics).}
The time evolution of a closed quantum system is governed by the
Schr\"{o}dinger equation
$i\hbar\,\partial_t\ket{\psi(t)} = H(t)\ket{\psi(t)}$,
where $H(t)$ is the Hermitian Hamiltonian of the system.
Equivalently, $\ket{\psi(t)} = U(t)\ket{\psi(0)}$ for some unitary $U(t)$.
\end{statement}

\paragraph{General solution via Dyson series.}
Formally integrating the Schr\"{o}dinger equation gives
\begin{align}
  \ket{\psi(t)} &= \mathcal{T}\exp\!\left(-\frac{i}{\hbar}\int_0^t H(t')\,dt'\right)\ket{\psi(0)},
  \label{eq:dyson}
\end{align}
where $\mathcal{T}$ denotes time-ordering.  For a \emph{time-independent}
Hamiltonian $H(t) = H$, the series resums exactly:
\begin{equation}
  \ket{\psi(t)} = e^{-iHt/\hbar}\ket{\psi(0)} \eqqcolon U(t)\ket{\psi(0)}.
  \label{eq:time-indep}
\end{equation}

\begin{proofsketch}
Split $[0,t]$ into $L$ equal intervals of width $\delta=t/L$.  For small $\delta$,
$U(\delta)\approx I - i H\delta/\hbar$.  The total evolution is
$U(t)\approx (I-iH\delta/\hbar)^L \xrightarrow{L\to\infty} e^{-iHt/\hbar}$
by the definition of the matrix exponential.  Hermiticity of $H$ gives
$U^\dagger = e^{iH^\dagger t/\hbar} = e^{iHt/\hbar}$, hence
$U^\dagger U = e^{iHt/\hbar}e^{-iHt/\hbar} = I$, confirming unitarity. $\lozenge$
\end{proofsketch}

\paragraph{Spectral form and inner-product preservation.}
Since $H$ is Hermitian, its spectral decomposition reads
$H = \sum_k E_k \ket{E_k}\bra{E_k}$, giving
\begin{equation}
  U(t) = \sum_k e^{-iE_k t/\hbar} \ket{E_k}\bra{E_k}.
\end{equation}
Because $|e^{-iE_k t/\hbar}|=1$, the evolution merely rotates phases; it preserves
the norm and inner product:
$\braket{\psi(t)|\phi(t)} = \bra{\psi(0)}U^\dagger U\ket{\phi(0)} = \braket{\psi(0)|\phi(0)}$.

\paragraph{Circuit model.}
In quantum computing we set $\hbar = 1$.  A quantum circuit decomposes
$U = U_L \cdots U_2 U_1$ where each $U_k$ acts on one or two qubits.  Because a product
of unitaries is unitary ($(VU)^\dagger(VU)=U^\dagger V^\dagger V U= I$),
every circuit realizes a valid unitary.

\subsection{Postulate III: Measurement}
\label{ssec:measurement}

\begin{statement}
\textbf{Postulate III (Measurement).}
Quantum measurements are described by a collection $\{M_m\}$ of measurement
operators satisfying the completeness relation $\sum_m M_m^\dagger M_m = \mathbf{I}$.
If the pre-measurement state is $\ket{\psi}$, the outcome $m$ occurs with
probability $p(m) = \bra{\psi}M_m^\dagger M_m\ket{\psi}$, and the
post-measurement state is $M_m\ket{\psi}/\sqrt{p(m)}$.
\end{statement}

\paragraph{Projective measurements.}
A \emph{projective} (von Neumann) measurement is associated with a Hermitian
observable $\mathcal{O} = \sum_m \mu_m P_m$, where $\mu_m$ are real eigenvalues
and $P_m = \ket{\mu_m}\bra{\mu_m}$ are orthogonal projectors:
$P_m P_{m'} = \delta_{mm'} P_m$, $\sum_m P_m = \mathbf{I}$.
Setting $M_m = P_m$:
\begin{equation}
  p(\mu_m) = \braket{\psi|P_m|\psi}, \qquad
  \ket{\psi'} = \frac{P_m\ket{\psi}}{\sqrt{p(\mu_m)}}.
\end{equation}

\paragraph{Expectation value.}
\begin{equation}
  \langle\mathcal{O}\rangle = \sum_m \mu_m p(\mu_m)
    = \bra{\psi}\mathcal{O}\ket{\psi}
    = \operatorname{Tr}(\mathcal{O}\,\ket{\psi}\bra{\psi}).
  \label{eq:expectation}
\end{equation}
The second equality follows from
$\sum_m \mu_m\braket{\psi|P_m|\psi} = \bra{\psi}(\sum_m \mu_m P_m)\ket{\psi}$.
The third uses the cyclic property of the trace.

\paragraph{Computational basis measurement.}
Setting $M_m = \ket{m}\bra{m}$ for $m \in \{0,1\}^n$ and
$\ket{\psi} = \sum_x \alpha_x\ket{x}$, outcome $m$ occurs with probability
$|\alpha_m|^2$, and the post-measurement state collapses to $\ket{m}$.

% ============================================================
\section{Multi-Qubit Systems}
\label{sec:multiqubit}
% ============================================================

\subsection{Composite Systems and Tensor Products}
\label{ssec:composite}

\begin{statement}
\textbf{Postulate IV (Composite Systems).}
The state space of a composite physical system is the tensor product
$\mathcal{H}_1 \otimes \cdots \otimes \mathcal{H}_n$ of the component Hilbert spaces.
If subsystem $k$ is in state $\ket{\psi_k}$, the joint state is
$\ket{\psi_1}\otimes\cdots\otimes\ket{\psi_n}$.
\end{statement}

For two qubits with $\ket{\psi_1} = \alpha_1\ket{0}+\beta_1\ket{1}$ and
$\ket{\psi_2} = \alpha_2\ket{0}+\beta_2\ket{1}$, the joint state is
\begin{align}
  \ket{\psi_1}\otimes\ket{\psi_2} &=
    \alpha_1\alpha_2\ket{00} + \alpha_1\beta_2\ket{01} \nonumber \\
    &\quad + \beta_1\alpha_2\ket{10} + \beta_1\beta_2\ket{11},
  \label{eq:product}
\end{align}
which coincides with the Kronecker product of the column vectors.  An $n$-qubit
Hilbert space has dimension $2^n$, growing exponentially---the root of both the
power and the classical-simulation difficulty of quantum computation.

\subsection{Quantum Entanglement}
\label{ssec:entanglement}

A bipartite state $\ket{\Psi} \in \mathcal{H}_A \otimes \mathcal{H}_B$ is
\emph{separable} if $\ket{\Psi} = \ket{\psi_A}\otimes\ket{\psi_B}$;
otherwise it is \emph{entangled}.  The four \emph{Bell states},
\begin{align}
  \ket{\Phi^{\pm}} &= \tfrac{1}{\sqrt{2}}(\ket{00} \pm \ket{11}), \\
  \ket{\Psi^{\pm}} &= \tfrac{1}{\sqrt{2}}(\ket{01} \pm \ket{10}),
\end{align}
are maximally entangled two-qubit states and form an orthonormal basis for
$\mathbb{C}^4$ (the Bell basis).

\paragraph{Schmidt decomposition and separability criterion.}
Any $\ket{\Psi} \in \mathcal{H}_A \otimes \mathcal{H}_B$ admits the
\emph{Schmidt decomposition}
\begin{equation}
  \ket{\Psi} = \sum_k \sqrt{\lambda_k}\, \ket{a_k}\ket{b_k},
  \quad \lambda_k \ge 0,\quad \sum_k \lambda_k = 1,
  \label{eq:schmidt}
\end{equation}
where $\{\ket{a_k}\}$ and $\{\ket{b_k}\}$ are orthonormal.  The Schmidt rank
(number of non-zero $\lambda_k$) characterizes entanglement: a state is separable
iff its Schmidt rank is 1.

\begin{proofsketch}
Represent $\ket{\Psi}$ as a matrix $M_{ij} = \alpha_{ij}$ via
$\ket{\Psi}=\sum_{i,j}\alpha_{ij}\ket{a_i}\ket{b_j}$.  Applying the singular
value decomposition (SVD) $M = U\Sigma V^\dagger$ gives
$\ket{\Psi}=\sum_k \sigma_k \ket{\tilde{a}_k}\ket{\tilde{b}_k}$ where
$\ket{\tilde{a}_k}=\sum_i U_{ik}\ket{a_i}$ and
$\ket{\tilde{b}_k}=\sum_j V^*_{jk}\ket{b_j}$ are new orthonormal sets.
Setting $\sqrt{\lambda_k}=\sigma_k$ and normalizing completes the proof. $\lozenge$
\end{proofsketch}

\subsection{Partial Measurement}
\label{ssec:partial}

Consider a two-qubit state $\ket{\Psi} = \sum_{i,j} \alpha_{ij}\ket{ij}$ with
$\sum_{i,j}|\alpha_{ij}|^2=1$.  Measuring only the first qubit with
$M_0 = \ket{0}\bra{0}\otimes\mathbf{I}$ and $M_1 = \ket{1}\bra{1}\otimes\mathbf{I}$:
\begin{itemize}
  \item Outcome $0$ with probability
        $p(0) = |\alpha_{00}|^2 + |\alpha_{01}|^2$;
        post-state $\ket{0}\otimes\frac{\alpha_{00}\ket{0}+\alpha_{01}\ket{1}}{\sqrt{p(0)}}$.
  \item Outcome $1$ with probability
        $p(1) = |\alpha_{10}|^2 + |\alpha_{11}|^2$;
        post-state $\ket{1}\otimes\frac{\alpha_{10}\ket{0}+\alpha_{11}\ket{1}}{\sqrt{p(1)}}$.
\end{itemize}
Notably, if $\ket{\Psi}=\ket{\Phi^+}$, measuring the first qubit and obtaining
$0$ instantly collapses the second qubit to $\ket{0}$, regardless of distance.
This is \emph{not} superluminal signaling: the outcome probabilities are always $1/2$.

\subsection{Pauli Matrices}
\label{ssec:pauli}

The Pauli matrices
\begin{equation}
  X = \begin{pmatrix}0&1\\1&0\end{pmatrix}, \quad
  Y = \begin{pmatrix}0&-i\\i&0\end{pmatrix}, \quad
  Z = \begin{pmatrix}1&0\\0&-1\end{pmatrix}
  \label{eq:pauli}
\end{equation}
together with $I\equiv\sigma_0$ satisfy $P^2=I$, $[X,Y]=2iZ$ (and cyclic),
and $\{P,Q\}=2\delta_{PQ}I$.  The set $\mathcal{P}_n = \{I,X,Y,Z\}^{\otimes n}$
(with phases $\{\pm 1,\pm i\}$) forms the $n$-qubit Pauli group and provides a
basis for $2^n\times 2^n$ Hermitian matrices.  This makes Pauli operators the
natural language for quantum error correction and VQA cost functions.

\subsection{Density Matrix Formalism}
\label{ssec:density}

When the state of a system is only known probabilistically---system is in
$\ket{\psi_k}$ with probability $p_k$---the \emph{density operator}
\begin{equation}
  \rho = \sum_k p_k \ket{\psi_k}\bra{\psi_k}
  \label{eq:density}
\end{equation}
provides a complete description.  Key properties:
(i) $\rho \ge 0$,
(ii) $\operatorname{Tr}(\rho) = 1$,
(iii) $\operatorname{Tr}(\rho^2) \le 1$, with equality iff $\rho$ is pure.
Under unitary evolution, $\rho \to U\rho U^\dagger$.  Under measurement $\{M_m\}$:
\begin{equation}
  p(m) = \operatorname{Tr}(M_m^\dagger M_m \rho), \qquad
  \rho' = \frac{M_m \rho M_m^\dagger}{p(m)}.
\end{equation}
The expectation value generalizes Eq.~\eqref{eq:expectation} to
$\langle\mathcal{O}\rangle = \operatorname{Tr}(\mathcal{O}\rho)$.

\subsection{Reduced Density Matrix and Partial Trace}
\label{ssec:reduced-density}

A central problem in quantum information is describing a subsystem $A$ of a composite system $AB$ when we lack access to system $B$. If the joint system is in state $\rho_{AB}$, the \emph{reduced density matrix} $\rho_A$ is defined via the \emph{partial trace} over $B$:
\begin{equation}
  \rho_A \equiv \operatorname{Tr}_B(\rho_{AB}) = \sum_j \bra{j}_B \rho_{AB} \ket{j}_B,
  \label{eq:partial-trace}
\end{equation}
where $\{\ket{j}_B\}$ is any orthonormal basis for $\mathcal{H}_B$.

\paragraph{Proof of Physical Consistency.}
The reduced density matrix $\rho_A$ is the unique operator that correctly predicts the expectation values of all local observables $O_A$ on system $A$.
\begin{proofsketch}
Let $O = O_A \otimes I_B$ be an observable acting only on system $A$. Its expectation value in state $\rho_{AB}$ is:
\begin{align}
  \langle O_A \otimes I_B \rangle &= \operatorname{Tr}\bigl((O_A \otimes I_B) \rho_{AB}\bigr) \nonumber \\
  &= \sum_i \sum_j \bra{i}_A \bra{j}_B (O_A \otimes I_B) \rho_{AB} \ket{i}_A \ket{j}_B \nonumber \\
  &= \sum_i \bra{i}_A O_A \underbrace{\left( \sum_j \bra{j}_B \rho_{AB} \ket{j}_B \right)}_{\operatorname{Tr}_B(\rho_{AB})} \ket{i}_A \nonumber \\
  &= \sum_i \bra{i}_A O_A \rho_A \ket{i}_A = \operatorname{Tr}(O_A \rho_A). \nonumber
\end{align}
This identity proves that $\rho_A$ contains all information necessary to describe local experiments on $A$ \cite{wilde2017quantum}. $\lozenge$
\end{proofsketch}

\paragraph{Entanglement and Mixedness.}
Consider the maximally entangled Bell state $\ket{\Phi^+} = \tfrac{1}{\sqrt{2}}(\ket{00}+\ket{11})$. Its joint density matrix is $\rho_{AB} = \ket{\Phi^+}\bra{\Phi^+}$. Computing the partial trace over $B$:
\begin{align}
  \rho_A &= \bra{0}_B \rho_{AB} \ket{0}_B + \bra{1}_B \rho_{AB} \ket{1}_B \nonumber \\
  &= \tfrac{1}{2}(\ket{0}\bra{0} + \ket{1}\bra{1}) = \tfrac{1}{2}\mathbf{I}.
\end{align}
While the joint state $\rho_{AB}$ is pure ($\text{Tr}(\rho_{AB}^2)=1$), the reduced state $\rho_A$ is \emph{maximally mixed} ($\text{Tr}(\rho_A^2)=1/2$). This conversion of global entanglement into local classical uncertainty (entropy) is a hallmark of quantum correlations \cite{nielsen2010quantum}.

% ============================================================
\section{Quantum Circuit Model}
\label{sec:circuits}
% ============================================================

\subsection{Motivation: Classical Circuits}
\label{ssec:classical-circuits}

Classical digital computation is described by Boolean circuits: directed acyclic
graphs of \emph{gates} with \emph{wires} carrying bits.  The set
$\{\mathrm{NAND}, \mathrm{FANOUT}\}$ is universal.  However, $\mathrm{NAND}$ is
\emph{irreversible} (many-to-one), which by Landauer's principle implies
thermodynamic energy dissipation (Sec.~\ref{ssec:landauer}).

\subsection{Single-Qubit Gates}
\label{ssec:single-qubit}

Single-qubit gates are $2\times 2$ unitary matrices.  The most important examples:
\begin{align}
  H &= \tfrac{1}{\sqrt{2}}\begin{pmatrix}1&1\\1&-1\end{pmatrix}, \quad
  X = \begin{pmatrix}0&1\\1&0\end{pmatrix}, \quad
  Z = \begin{pmatrix}1&0\\0&-1\end{pmatrix}, \\
  S &= \begin{pmatrix}1&0\\0&i\end{pmatrix}, \quad
  T  = \begin{pmatrix}1&0\\0&e^{i\pi/4}\end{pmatrix}, \quad T^2=S,\;S^2=Z.
\end{align}
Rotation gates about Bloch-sphere axes:
\begin{equation}
  R_{\hat{n}}(\theta) \coloneqq e^{-i\theta\hat{n}\cdot\vec{\sigma}/2}
    = \cos\!\tfrac{\theta}{2}\,I - i\sin\!\tfrac{\theta}{2}\,\hat{n}\cdot\vec{\sigma},
  \label{eq:rotation}
\end{equation}
where $\vec{\sigma} = (X,Y,Z)$.  Geometrically, $R_{\hat{n}}(\theta)$ rotates the
Bloch-sphere vector by angle $\theta$ about the $\hat{n}$-axis.

\paragraph{Euler decomposition.}
Every single-qubit unitary $U \in SU(2)$ can be written as
\begin{equation}
  U = e^{i\alpha}R_z(\beta)\,R_y(\gamma)\,R_z(\delta)
  \label{eq:euler}
\end{equation}
for some real $\alpha,\beta,\gamma,\delta$.

\begin{proofsketch}
Since $U$ is $2\times2$ unitary and $\det U=1$ (up to global phase), write
$U = \bigl(\begin{smallmatrix}a & -b^* \\ b & a^*\end{smallmatrix}\bigr)$ with
$|a|^2+|b|^2=1$.  Set $a = e^{i(\alpha-\beta/2-\delta/2)}\cos(\gamma/2)$
and $b = e^{i(\alpha+\beta/2-\delta/2)}\sin(\gamma/2)$.
Matching this to $R_z(\beta)R_y(\gamma)R_z(\delta)$ computed explicitly via
Eq.~\eqref{eq:rotation} confirms the decomposition for all $U \in SU(2)$. $\lozenge$
\end{proofsketch}

A hardware-convenient parameterization for superconducting qubits is
\begin{equation}
  U(\theta,\phi,\lambda) = R_z(\phi)\,R_x\!\left(-\tfrac{\pi}{2}\right)\,R_z(\theta)\,R_x\!\left(\tfrac{\pi}{2}\right)\,R_z(\lambda).
  \label{eq:ibm-u3}
\end{equation}

\subsection{Two-Qubit Gates and Entanglement Generation}
\label{ssec:two-qubit}

Product operators $U_A \otimes U_B$ cannot generate entanglement.  The canonical
entangling gates are:
\begin{align}
  CX &= \ket{0}\bra{0}\otimes I + \ket{1}\bra{1}\otimes X, \\
  CZ &= \ket{0}\bra{0}\otimes I + \ket{1}\bra{1}\otimes Z.
\end{align}
$CX$ flips the target conditioned on the control being $\ket{1}$.  Applying $CX$
to $(H\ket{0})\otimes\ket{0} = \tfrac{1}{\sqrt{2}}(\ket{0}+\ket{1})\otimes\ket{0}$:
\begin{equation}
  CX\cdot\tfrac{1}{\sqrt{2}}(\ket{00}+\ket{10})
  = \tfrac{1}{\sqrt{2}}(\ket{00}+\ket{11}) = \ket{\Phi^+},
\end{equation}
demonstrating Bell-state generation from a product state.

\subsection{Universality of Quantum Gates}
\label{ssec:universality}

\begin{statement}
\textbf{Theorem (Universality).}
A gate set $\mathcal{G}$ is \emph{universal for quantum computation} if any
$n$-qubit unitary can be approximated to precision $\varepsilon$ by a circuit of
size $O\!\left(\operatorname{poly}\!\left(\log\tfrac{1}{\varepsilon}\right)\right)$ from $\mathcal{G}$.
The set $\{H, T, CX\}$ is universal \cite{nielsen2010quantum}.
\end{statement}

\begin{proofsketch}
The argument proceeds in two steps.
(1) \textit{Single-qubit universality}: $H$ and $T$ together generate a dense
subgroup of $SU(2)$---this can be shown by computing $HTH$ and checking that the
resulting rotations are irrational multiples of $\pi$, ensuring density.
(2) \textit{Multi-qubit reduction}: any $n$-qubit unitary can be decomposed into
two-qubit unitaries via a sequence of controlled-$U$ gates, each decomposable into
$CX$ and single-qubit gates (Gray-code argument).  The Solovay-Kitaev theorem then
guarantees that approximation within $\varepsilon$ requires only
$O(\operatorname{poly}\log(1/\varepsilon))$ gates. $\lozenge$
\end{proofsketch}

\subsection{No-Cloning Theorem}
\label{ssec:no-cloning}

\begin{statement}
\textbf{Theorem (No-Cloning, Wootters \& Zurek 1982).}
There is no unitary operation $U$ such that
$U\ket{\psi}\ket{0} = \ket{\psi}\ket{\psi}$ for all $\ket{\psi}$
\cite{wootters1982single}.
\end{statement}

\begin{proofsketch}[Proof]
Suppose such $U$ exists and let $\ket{\psi},\ket{\phi}$ be arbitrary.  Then
$U\ket{\psi}\ket{0} = \ket{\psi}\ket{\psi}$ and
$U\ket{\phi}\ket{0} = \ket{\phi}\ket{\phi}$.
Taking the inner product of both equations and using unitarity ($U^\dagger U=I$):
\begin{equation}
  \braket{\psi|\phi}\underbrace{\braket{0|0}}_{=1}
  = \braket{\psi|\phi}\braket{\psi|\phi}
  = \braket{\psi|\phi}^2.
\end{equation}
Hence $\braket{\psi|\phi}(\braket{\psi|\phi}-1) = 0$, forcing
$\braket{\psi|\phi} \in \{0,1\}$.  Any two states that $U$ clones must be either
identical or orthogonal, contradicting the assumption that $U$ clones
\emph{all} states.
\end{proofsketch}

\paragraph{Implication for QML.}
No-cloning forbids direct copying of unknown quantum data, restricting
classical-style data augmentation.  It simultaneously motivates quantum-native
protocols (superdense coding, teleportation) that exploit entanglement instead
of copying.

\subsection{Basic Quantum Protocols}
\label{ssec:protocols}

\paragraph{Superdense coding.}
By sharing the Bell pair $\ket{\Phi^+}=\tfrac{1}{\sqrt{2}}(\ket{00}+\ket{11})$
in advance, Alice can transmit \emph{2 classical bits} to Bob by sending only a
\emph{single qubit} \cite{bennett1992communication}.  By applying a local unitary
$X^a Z^b$ ($a,b \in \{0,1\}$) to her qubit, Alice can transform the globally shared
entangled state into any of the four orthogonal Bell states.  Bob then performs
a Bell-basis measurement to distinguish these states with 100\% certainty.

\begin{figure}[h]
\centering
\begin{quantikz}
  \lstick[2]{$\ket{\beta_{00}}$} & \gate{Z^b} & \gate{X^a} & \qw & \qw & \targ{} \gategroup[2,steps=3,style={dashed, rounded corners, draw=red, inner sep=2pt}, background]{Bell measurement} & \qw & \meter{} & \rstick{a} \\
  & \qw      & \qw  & \qw & \qw & \ctrl{-1}  & \gate{H} & \meter{} & \rstick{b}
\end{quantikz}
\caption{Superdense coding circuit.  Alice applies local operators $Z^b$ and $X^a$ to her
half of a shared Bell pair ($q_A$).  She then sends $q_A$ to Bob, who performs a 
Bell-basis measurement (CNOT and $H$) to recover both bits $a, b$.}
\label{fig:superdense}
\end{figure}

\begin{table*}[t]
\centering
\caption{Alice's encoding and resulting Bell states for Superdense Coding. The bit $b$ (determined by the wire with the $H$ gate) controls the $Z$ gate.}
\label{tab:superdense}
\begin{tabular}{@{}cccc@{}}
\toprule
Requested Bits ($ba$) & Alice's Gate Logic ($X^a Z^b$) & Resulting Bell State (on wire A) & Bob's Measurement Result ($ba$) \\ \midrule
00 & $I$  & $\ket{\Phi^+} = \tfrac{1}{\sqrt{2}}(\ket{00}+\ket{11})$ & 00 \\
01 & $X$  & $\ket{\Psi^+} = \tfrac{1}{\sqrt{2}}(\ket{10}+\ket{01})$ & 01 \\
10 & $Z$  & $\ket{\Phi^-} = \tfrac{1}{\sqrt{2}}(\ket{00}-\ket{11})$ & 10 \\
11 & $XZ$ & $\ket{\Psi^-} = \tfrac{1}{\sqrt{2}}(\ket{10}-\ket{01})$ & 11 \\ \bottomrule
\end{tabular}
\end{table*}

The protocol succeeds because the set $\{ (E_i \otimes I) \ket{\Phi^+} \}$ forms an 
orthonormal basis for the four-dimensional Hilbert space $(\mathbb{C}^2)^{\otimes 2}$.
Alice effectively picks which basis vector the system resides in, and Bob performs
the change-of-basis back to the computational basis for sampling.


\paragraph{Quantum teleportation.}
Alice transmits an unknown qubit $\ket{\psi}=\alpha\ket{0}+\beta\ket{1}$ to Bob
using only a shared Bell pair and two classical bits \cite{bennett1993teleporting}.
The circuit (Fig.~\ref{fig:teleport}) evolves as follows.  The three-qubit
initial state is
\begin{equation}
  \ket{\psi}\otimes\ket{\Phi^+}
  = \tfrac{1}{\sqrt{2}}\bigl(\alpha\ket{0}+\beta\ket{1}\bigr)(\ket{00}+\ket{11}).
\end{equation}
After Alice applies $CX$ then $H$ on her qubits, the state becomes
\begin{align}
  \tfrac{1}{2}\bigl[&\ket{00}(\alpha\ket{0}+\beta\ket{1})
                    +\ket{01}(\alpha\ket{1}+\beta\ket{0}) \nonumber\\
                   +&\ket{10}(\alpha\ket{0}-\beta\ket{1})
                    +\ket{11}(\alpha\ket{1}-\beta\ket{0})\bigr].
\end{align}
Alice measures qubits 1 and 2 and sends outcomes $(a,b)$ classically.  Bob applies
$Z^a X^b$ to recover $\alpha\ket{0}+\beta\ket{1}=\ket{\psi}$ exactly.

\begin{figure}[h]
\centering
\begin{quantikz}
  \lstick{A $\ket{\psi}$} & \ctrl{1} & \gate{H} & \qw & \meter{}  & \cwbend{2}  \\
  \lstick[2]{$\ket{\beta_{00}}$} & \targ{}  & \qw & \meter{} & \cwbend{1} \\
  & \qw      & \qw      & \qw      & \gate{X^b} & \gate{Z^a} & \rstick{$\ket{\psi}$}
\end{quantikz}
\caption{Quantum teleportation circuit.  Alice performs a Bell measurement on
her state $A$ and her half of a shared Bell pair $A$.  The outcomes $a$ and $b$ 
are used to apply gates $X^b$ and $Z^a$ to Bob's qubit $B$, recovering $\ket{\psi}$.}
\label{fig:teleport}
\end{figure}

\begin{table*}[t]
\centering
\caption{Alice's measurement outcomes and Bob's corresponding recovery gates. The bit $a$ (from the wire with $H$) controls the $Z$ gate.}
\label{tab:teleport}
\begin{tabular}{@{}cccc@{}}
\toprule
Outcome ($ba$) & Alice's Projector & Bob's Received State & Correction Logic ($Z^a X^b$) \\ \midrule
00 & $\ket{\Phi^+}\bra{\Phi^+}$ & $\alpha\ket{0}+\beta\ket{1}$ & $I$ \\
10 & $\ket{\Psi^+}\bra{\Psi^+}$ & $\alpha\ket{1}+\beta\ket{0}$ & $X$ \\
01 & $\ket{\Phi^-}\bra{\Phi^-}$ & $\alpha\ket{0}-\beta\ket{1}$ & $Z$ \\
11 & $\ket{\Psi^-}\bra{\Psi^-}$ & $\alpha\ket{1}-\beta\ket{0}$ & $ZX$ \\ \bottomrule
\end{tabular}
\end{table*}

Importantly, the "teleportation" is not instantaneous. Although the system 
collapses into one of the four states in Table \ref{tab:teleport} immediately 
upon Alice's measurement, Bob cannot know which gate to apply without the two 
classical bits.  This preserves causality and satisfies the 
\emph{no-communication theorem}: entanglement alone cannot transmit 
information faster than light.


% ============================================================
\section{Reversible Computation and Complexity}
\label{sec:reversibility}
% ============================================================

\subsection{Landauer's Principle}
\label{ssec:landauer}

\begin{statement}
\textbf{Landauer's Principle (1961).}
Erasing one bit of information irreversibly dissipates at least $k_B T \ln 2$
of energy as heat, where $k_B$ is Boltzmann's constant and $T$ is the
temperature of the environment \cite{landauer1961irreversibility}.
\end{statement}

Since NAND (and most classical gates) are many-to-one, they are thermodynamically
irreversible and incur this minimum energy cost.  In contrast, unitary quantum
evolution is always reversible ($U^{-1} = U^\dagger$), satisfying Landauer's
bound with equality in principle.

\subsection{Toffoli Gate: Universal Reversible Classical Gate}
\label{ssec:toffoli}

The Toffoli (CCX) gate maps $(a, b, c) \mapsto (a, b, c \oplus ab)$.
It is universal for reversible classical computation: both NAND and FANOUT can be
simulated using Toffoli gates with ancilla bits.  As a quantum gate on three qubits,
it is a $8\times 8$ permutation matrix.

\paragraph{Quantum oracle for arbitrary functions.}
Any $f:\{0,1\}^n \to \{0,1\}^m$ can be embedded into the reversible unitary
\begin{equation}
  U_f : \ket{x}\ket{y} \mapsto \ket{x}\ket{y \oplus f(x)},
  \label{eq:oracle}
\end{equation}
where $\oplus$ denotes bitwise XOR.  Applying $U_f$ to a uniform superposition:
\begin{equation}
  U_f\!\left(\frac{1}{\sqrt{2^n}}\sum_x \ket{x}\right)\!\ket{0}
  = \frac{1}{\sqrt{2^n}}\sum_x \ket{x}\ket{f(x)},
\end{equation}
evaluating $f$ on all $2^n$ inputs simultaneously---the key mechanism underlying
quantum speedups via interference.

\subsection{Hadamard Test and Swap Test}
\label{ssec:hadamard-test}

\paragraph{Hadamard test.}
Given a unitary $U$ and a state $\ket{\psi}$, the Hadamard test
(Fig.~\ref{fig:hadamard}) estimates $\langle\psi|U|\psi\rangle \in \mathbb{C}$.
A full step-by-step derivation follows.

\begin{proofsketch}[Complete derivation of Hadamard test]
\textit{Step 1.} Initialize ancilla and system:
$\ket{0}\ket{\psi}$.

\textit{Step 2.} Apply $H$ to ancilla:
$\tfrac{1}{\sqrt{2}}(\ket{0}+\ket{1})\ket{\psi}$.

\textit{Step 3.} Apply controlled-$U$ (if ancilla is $\ket{1}$, apply $U$ to system):
\begin{equation}
  \tfrac{1}{\sqrt{2}}\bigl(\ket{0}\ket{\psi} + \ket{1}U\ket{\psi}\bigr).
\end{equation}

\textit{Step 4.} Apply final $H$ to ancilla:
\begin{align}
  &\tfrac{1}{2}\bigl[(\ket{0}+\ket{1})\ket{\psi}
   + (\ket{0}-\ket{1})U\ket{\psi}\bigr] \nonumber\\
  &= \tfrac{1}{2}\ket{0}(\ket{\psi}+U\ket{\psi})
   + \tfrac{1}{2}\ket{1}(\ket{\psi}-U\ket{\psi}).
\end{align}

\textit{Step 5.} Measure ancilla.  The outcome probabilities are:
\begin{align}
  P(0) &= \tfrac{1}{4}\|\ket{\psi}+U\ket{\psi}\|^2
       = \tfrac{1}{4}(2 + 2\operatorname{Re}\braket{\psi|U|\psi}), \\
  P(1) &= \tfrac{1}{4}(2 - 2\operatorname{Re}\braket{\psi|U|\psi}).
\end{align}
Hence $P(0)-P(1) = \operatorname{Re}\braket{\psi|U|\psi}$.

\textit{Imaginary part.} Insert an $S^\dagger$ phase gate before the final
$H$; a parallel argument gives $P(0)-P(1) = \operatorname{Im}\braket{\psi|U|\psi}$.
Together, both runs fully specify $\braket{\psi|U|\psi} \in \mathbb{C}$. $\lozenge$
\end{proofsketch}

\begin{figure}[h]
\centering
\begin{quantikz}
  \lstick{$\ket{0}$} & \gate{H} & \ctrl{1} & \gate{H} & \meter{} & \rstick{$\sigma_z$} \\
  \lstick{$\ket{\psi}$} & \qw & \gate{U} & \qw & \qw & \qw
\end{quantikz}
\caption{Hadamard test circuit.  The ancilla qubit $\ket{0}$ controls the
unitary $U$ on the system register $\ket{\psi}$.  A final Hadamard on the
ancilla followed by $\sigma_z$ measurement yields
$P(0)-P(1)=\mathrm{Re}\langle\psi|U|\psi\rangle$.
Inserting $S^\dagger$ before the final $H$ yields the imaginary part.}
\label{fig:hadamard}
\end{figure}

\paragraph{Swap test.}
Setting $U=\mathrm{SWAP}$ and $\ket{\psi}=\ket{\psi_1}\ket{\psi_2}$, the
Hadamard test gives
\begin{equation}
  P(0)-P(1) = \langle\psi_1\psi_2|\mathrm{SWAP}|\psi_1\psi_2\rangle
  = |\langle\psi_1|\psi_2\rangle|^2.
\end{equation}
The identity $\bra{\psi_1\psi_2}\mathrm{SWAP}\ket{\psi_1\psi_2}
= \bra{\psi_1\psi_2}\ket{\psi_2\psi_1} = |\braket{\psi_1|\psi_2}|^2$
follows directly from the action of SWAP.

\begin{figure}[h]
\centering
\begin{quantikz}
  \lstick{$\ket{0}$} & \gate{H}  & \ctrl{1} & \gate{H} & \meter{} \\
  \lstick{$\ket{\psi_1}$} & \qw & \swap{1} & \qw & \qw \\
  \lstick{$\ket{\psi_2}$} & \qw & \targX{} & \qw & \qw
\end{quantikz}
\caption{Swap test circuit.  A controlled-SWAP between $\ket{\psi_1}$ and
$\ket{\psi_2}$ is mediated by the ancilla qubit.  The outcome probability
difference satisfies $P(0)-P(1) = |\langle\psi_1|\psi_2\rangle|^2$, providing
an efficient quantum estimate of state overlap used in kernel-based QML
\cite{schuld2019quantum}.}
\label{fig:swap}
\end{figure}


\subsection{Quantum Complexity Classes}
\label{ssec:complexity-classes}

\begin{itemize}
  \item \textbf{BQP} (Bounded-error Quantum Polynomial time): decision problems
    solvable by a polynomial-time quantum circuit with error $\le 1/3$.
    $\mathrm{BPP}\subseteq\mathrm{BQP}\subseteq\mathrm{PSPACE}$.
  \item \textbf{QMA} (Quantum Merlin-Arthur): the quantum analogue of NP;
    YES instances have a polynomial-size quantum witness verifiable by
    a BQP verifier.
\end{itemize}
Whether $\mathrm{BQP} \not\subseteq \mathrm{NP}$ or $\mathrm{BQP}\ne\mathrm{BPP}$
remains open \cite{bernstein1997quantum}.  The believed picture is
$\mathrm{BPP} \subsetneq \mathrm{BQP}$, with factoring as a candidate
separation (Shor's algorithm).

% ============================================================
\section{Towards Quantum Machine Learning}
\label{sec:qml}
% ============================================================

\subsection{Quantum Advantage in Machine Learning}
\label{ssec:qml-advantage}

Quantum machine learning (QML) investigates whether quantum computers can
accelerate or qualitatively improve machine learning tasks
\cite{biamonte2017quantum}.  Three criteria define meaningful quantum advantage:
\begin{enumerate}
  \item \textit{Physical feasibility} (with or without QEC),
  \item \textit{Computational advantage} over the best known classical algorithm,
  \item \textit{Industry relevance} for real-world ML tasks.
\end{enumerate}
Meeting all three simultaneously is an open problem.  Quantum speedups in
supervised learning often rely on quantum RAM (QRAM) for efficient data
loading---a hardware resource whose practical feasibility remains uncertain.

\subsection{Key QML Algorithms}
\label{ssec:qml-algorithms}

\paragraph{HHL Algorithm (Linear Systems).}
Harrow, Hassidim, and Lloyd \cite{harrow2009quantum} showed that the system
$A\mathbf{x} = \mathbf{b}$ (sparse $A\in\mathbb{C}^{N\times N}$) can be solved
quantumly in $O(\log N\cdot\kappa^2/\varepsilon)$ time, where $\kappa$ is the
condition number.  The algorithm uses quantum phase estimation to encode
eigenvalue inverses and outputs a quantum state $\ket{x}\propto A^{-1}\ket{b}$.

\paragraph{Quantum Principal Component Analysis.}
Lloyd, Mohseni, and Rebentrost \cite{lloyd2014quantum} estimated principal
components of a density matrix $\rho$ in $O(\log N)$ time---exponentially faster
than classical PCA under analogous QRAM assumptions.

\paragraph{Quantum Support Vector Machine.}
Rebentrost, Mohseni, and Lloyd \cite{rebentrost2014quantum} proposed a quantum
SVM achieving $O(\log N)$ classification via HHL-based dual optimization.
The swap test (Sec.~\ref{ssec:hadamard-test}) serves as the quantum kernel evaluator.

\paragraph{Variational Quantum Algorithms.}
VQAs use short parameterized circuits $U(\boldsymbol{\theta})$ and optimize
$\boldsymbol{\theta}$ classically to minimize
$C(\boldsymbol{\theta}) = \langle\psi(\boldsymbol{\theta})|
\mathcal{O}|\psi(\boldsymbol{\theta})\rangle$ \cite{cerezo2021variational}.
VQAs are the primary paradigm for NISQ-era QML \cite{preskill2018quantum}.

\subsection{Open Questions and Future Outlook}
\label{ssec:outlook}

Despite significant theoretical progress, fundamental questions remain:
\begin{itemize}
  \item \textit{Dequantization}: Can classical sampling-based algorithms
    match quantum speedups \cite{schuld2021machine}?
  \item \textit{Barren plateaus}: Do VQA gradients vanish exponentially
    with system size, making training intractable?
  \item \textit{End-to-end advantage}: Is a provable, QRAM-free quantum speedup
    achievable on relevant ML tasks?
\end{itemize}
Progress on these questions will determine whether QML transcends academic
interest to become a practical technology.

% ============================================================
\section{Conclusion}
\label{sec:conclusion}
% ============================================================

This report has developed the theoretical foundations of quantum machine learning
from first principles.  Beginning with the postulates of quantum mechanics and the
mathematical framework of Hilbert spaces and Dirac notation, we derived the
quantum circuit model, established universality, proved the no-cloning theorem,
and analyzed reversible computation and Landauer's principle.  Detailed circuit
derivations for the Hadamard and swap tests illuminated how quantum expectation
values and state overlaps are estimated---quantities central to kernel-based QML.
Finally, we surveyed landmark QML algorithms (HHL, QPCA, QSVM, VQA) and
highlighted the central open questions surrounding quantum advantage.

Subsequent sections of this report (to be added as the course progresses) will
cover quantum Fourier transform, Grover's search, quantum kernel methods, and
variational quantum algorithms in depth.

\bibliography{bibliography}

\end{document}% trigger rebuild
% trigger integrated build
% pipeline trigger
% permission fix trigger
% math fix trigger
% premium dashboard trigger
% final pipeline fix trigger
% npm fix trigger
% ultimate fix trigger
% pdf.js viewer trigger
% final restoration trigger
% absolute final trigger
% optimization trigger
% multi-job restoration trigger
% final simple trigger
